\documentclass[10pt]{article}
\usepackage[letterpaper,margin=1in]{geometry}
\usepackage[T1]{fontenc}
\usepackage{lmodern,microtype,amsmath,amssymb,amsthm}
\usepackage[hidelinks]{hyperref}
\newtheorem{theorem}{Theorem}
\newtheorem{lemma}{Lemma}
\newcommand{\rk}{\operatorname{rk}}
\newcommand{\Fp}{\mathbb F_p}
\title{A classical Wu-manifold answer\\ to a free-rank question}
\author{Alec Kriebel\thanks{Independent researcher. ORCID:
\href{https://orcid.org/0009-0001-9320-500X}{0009-0001-9320-500X}.}}
\date{16 October 2026\quad (version 1.0)}
\begin{document}
\maketitle
\begin{abstract}
The fixed Wu manifold \(W=SU(3)/SO(3)\) is a finite simply connected complex
with the rational homotopy type of \(S^5\). Its free toral rank and free
\(2\)-rank are zero, whereas its free \(p\)-rank is one for every odd prime.
Thus it answers the existence question printed in Hanke's 2010 Oberwolfach
contribution, using the actual-space free-action definitions there.
The proof assembles classical facts: an almost-free circle action supplies
every odd cyclic action, and a nonzero Stiefel--Whitney number excludes
every continuous free involution. We give the arguments explicitly,
including the arbitrary-action odd-primary upper bound. This is an
explanatory research note; historical priority remains unresolved and no
firstness claim is made.
\end{abstract}

\section{The printed question and the fixed witness}
For a finite connected complex \(X\), let
\[
 \rk_T(X)=\max\{r:(S^1)^r\text{ acts freely on }X\},\qquad
 \rk_p(X)=\max\{r:(\mathbb Z/p)^r\text{ acts freely on }X\}.
\]
Actions here are continuous and act on the actual fixed space.
In \cite[p.~1912]{H}, Hanke asks whether there is such an \(X\) with
abelian fundamental group acting trivially on its higher homotopy,
finite-dimensional total rational homotopy, and
\(\rk_p(X)>\rk_T(X)\) for infinitely many primes \(p\).
These ranks do not optimize over rationally or homotopy equivalent
replacements, and the torus action must be free rather than almost free.

\begin{theorem}\label{main}
Let \(W=SU(3)/SO(3)\), with the real-matrix embedding of \(SO(3)\).
It is a finite simply connected complex of rational homotopy type \(S^5\),
and
\[
 \rk_T(W)=\rk_2(W)=0,\qquad
 \rk_p(W)=1\quad\text{for every odd prime }p.
\]
Moreover, every cyclic group of odd order acts freely and smoothly on \(W\).
\end{theorem}

The same \(W\) works before the prime is chosen. This supplies an affirmative
answer to the literal existence question in \cite{H}.
Its rational homotopy Euler characteristic is \(-1\), so the asymptotic
bound \(\rk_p(X)\le-\chi^\pi(X)\) reported there is respected.
The circle used below is almost free. Consequently this example does
not give a strict gap against an almost-free or rational toral rank.

\section{Topology and a classical characteristic number}
The principal bundle
\[
 SO(3)\longrightarrow SU(3)\longrightarrow W
\]
makes \(W\) a closed smooth manifold of dimension \(8-3=5\).
Its finite smooth triangulation realizes its actual underlying space as
a finite complex. The homotopy sequence, using
\(\pi_1(SU(3))=\pi_2(SU(3))=0\), gives
\(\pi_1(W)=0\) and \(\pi_2(W)=\mathbb Z/2\).
Hurewicz, integral Poincare duality and universal coefficients give
\[
 H_i(W;\mathbb Z)=
 \begin{cases}
 \mathbb Z,&i=0,5,\\
 \mathbb Z/2,&i=2,\\
 0,&\text{otherwise}.
 \end{cases}
\]
In particular the degree-three torsion is in integral \emph{cohomology},
not homology. Rational Hurewicz provides a sphere map \(S^5\to W\)
nonzero on rational top homology; it is a rational homology equivalence.
Rational Whitehead applies to these simply connected spaces, proving
the asserted rational homotopy type and its finite total dimension.
The fundamental-group hypotheses of the question are automatic.

Write \(x\) and \(y\) for the nonzero classes in
\(H^2(W;\mathbb F_2)\) and \(H^3(W;\mathbb F_2)\), respectively.
The classical tangent characteristic-number calculation is
\begin{equation}\label{number}
 w_2(W)=x,\qquad w_3(W)=y,\qquad
 \langle xy,[W]\rangle=1.
\end{equation}
A modern explicit account is \cite[Lemma~4.41, Corollary~4.44,
Proposition~4.45]{DY}, which credits older work of Barden,
Landweber--Stong and Calabi. Here is a direct reconstruction of the part
needed.

The tangent isotropy representation of \(SO(3)\) is its conjugation action
on \(i\,\mathrm{Sym}_0(\mathbb R^3)\), by
\(\mathfrak{su}(3)=\mathfrak{so}(3)\oplus i\,\mathrm{Sym}_0(\mathbb R^3)\).
A coordinate rotation splits this representation into a fixed line and
two real rotation planes of weights \(1\) and \(2\). Its map
\(\pi_1(SO(3))\to\pi_1(SO(5))\) is therefore nonzero, since
\(1+2\equiv1\pmod2\).
The boundary \(\pi_2(W)\to\pi_1(SO(3))\) is an isomorphism.
Pulling the tangent bundle to a representative generator sphere gives
the nontrivial clutching bundle on \(S^2\), so its \(w_2\) is nonzero.
The sphere's integral Hurewicz class generates \(\mathbb Z/2\) and remains
nonzero after reduction mod \(2\). Thus it detects \(w_2(W)=x\).

For \(0\to\mathbb Z/2\to\mathbb Z/4\to\mathbb Z/2\to0\),
universal coefficients identify the degree-two coefficient reduction with
\[
 \operatorname{Hom}(\mathbb Z/2,\mathbb Z/4)
 \longrightarrow\operatorname{Hom}(\mathbb Z/2,\mathbb Z/2).
\]
Both possible values \(0,2\) reduce to zero. The connecting map \(Sq^1\)
is consequently nonzero on \(x\), giving \(Sq^1x=y\).
Orientation and \(Sq^1w_2=w_1w_2+w_3\) give \(w_3=y\).
The complementary one-dimensional mod-\(2\) cohomology groups pair
nondegenerately, so \(\langle xy,[W]\rangle=1\).
These are tangent classes, not merely the characteristic classes of the
principal \(SO(3)\)-bundle.

\section{Odd cyclic actions and the obstruction at two}
For \(z\in S^1\), set
\[
 D(z)=\operatorname{diag}(z,z,z^{-2})\in SU(3),\qquad
 z\cdot[g]=[D(z)g].
\]
If \(z\) fixes a coset, then \(g^{-1}D(z)g\in SO(3)\).
Every determinant-one real orthogonal matrix of odd size has eigenvalue
\(1\). Thus \(z=1\) or \(z^{-2}=1\), so every stabilizer is contained in
\(\{1,-1\}\).
Every finite subgroup of odd order therefore acts freely. This includes
order \(3\): its nonidentity scalar matrices have no eigenvalue \(1\).
The full circle is not free, since
\(D(-1)=\operatorname{diag}(-1,-1,1)\in SO(3)\) fixes the identity coset.
This is the \((1,1,-2)\) instance of the classical circle-isotropy
calculation in \cite[Proposition~4.2]{Y}.

\begin{lemma}\label{cover}
A closed topological manifold with a nonzero tangent Stiefel--Whitney
number admits no free continuous involution.
\end{lemma}
\begin{proof}
A free involution has disjoint translated coordinate neighborhoods.
Its quotient \(N\) is a closed topological manifold and
\(q:M\to N\) an ordinary two-sheeted covering.
Local diagonal neighborhoods identify the topological tangent
microbundle of \(M\) with \(q^*\tau_N\).
Naturality of the mod-\(2\) Thom-class and Steenrod-square definition of
Stiefel--Whitney classes gives \(w_i(M)=q^*w_i(N)\).
With mod-\(2\) fundamental classes, \(q_*[M]=2[N]=0\).
Hence, for every partition \(i_1+\cdots+i_k=\dim M\),
\[
 \langle w_{i_1}(M)\cdots w_{i_k}(M),[M]\rangle
 =\langle w_{i_1}(N)\cdots w_{i_k}(N),q_*[M]\rangle=0,
\]
a contradiction. Smoothness, local linearity and orientability of \(N\)
are unnecessary.
\end{proof}

Apply Lemma~\ref{cover} to \eqref{number}. There is no free continuous
involution on \(W\), so \(\rk_2(W)=0\).
Every positive-dimensional torus contains an involution, and restriction
preserves freeness, proving \(\rk_T(W)=0\).
The nonbounding obstruction to a Wu involution is already explicitly
discussed in \cite[Theorem~9.1, v1; Example~2.25, v2]{KL}.
The covering argument here supplies the full continuous-action scope
without importing the admissibility condition of the surrounding
chromatic fixed-point theorem.
It uses only a finite-group quotient and makes no assertion about orbit
manifolds of arbitrary continuous circle actions.

\section{The exact rank at every odd prime}
For odd \(p\), the integral homology calculation makes \(W\) an
\(\Fp\)-homology five-sphere.
Suppose \(E=(\mathbb Z/p)^2\) acted freely and continuously.
Its action on the two nonzero fiber-cohomology lines is trivial, because
a \(p\)-group maps trivially to \(\Fp^\times\).
The Borel fibration has two-row Serre spectral sequence
\[
 E_2^{i,j}=H^i(BE;\Fp)\otimes H^j(W;\Fp),\qquad
 R=H^*(BE;\Fp)=\Fp[t_1,t_2]\otimes\Lambda(s_1,s_2),
\]
where \(|t_i|=2\) and \(|s_i|=1\).
Only \(d_6\) can be nonzero. Its bottom-row image is the principal ideal
\((f)\), where \(f=d_6(u)\) for the degree-five fiber generator.
Thus \(E_\infty^{*,0}=R/(f)\).

Killing the exterior generators gives a surjection
\[
 R/(f)\longrightarrow\Fp[t_1,t_2]/(\bar f).
\]
If \(\bar f=0\), the target is infinite-dimensional.
Otherwise \(\bar f\) is a nonzero homogeneous binary cubic.
In ordinary polynomial degree \(m\ge3\), multiplication by it is
injective from a space of dimension \(m-2\) into one of dimension \(m+1\);
the quotient has dimension \(3\).
Thus bottom-row classes survive in arbitrarily high total degrees.
This argument permits exterior terms in \(f\), and does not assume that
\(f\) is nonzero or a non-zero-divisor in \(R\).
The first-quadrant sequence has a finite filtration in each total degree;
extensions cannot erase its nonzero associated-graded pieces.

Freeness makes \(W/E\) a closed topological five-manifold and identifies
the Borel construction with it up to homotopy.
Its cohomology vanishes above degree five, a contradiction.
Higher-rank elementary abelian actions restrict to rank two and are
excluded as well. Together with the cyclic construction this proves
\(\rk_p(W)=1\) for every odd \(p\), completing Theorem~\ref{main}.
The universal conclusions follow from the proofs, not from sampling primes.

\section{Attribution, verification and status}
Earlier actual free finite-versus-circle gaps must be distinguished from
the extra hypotheses in Hanke's question.
Mundet i Riera \cite[Theorem~1.11, Section~13]{M}, already in the first
2021 preprint as Theorem~1.14 and Section~15, constructs regular self-covers
of every odd degree of the fixed Cappell--Weinberger--Yan manifolds.
These are circle-cover pullbacks, so every odd prime gives a genuinely
free cyclic deck action; the manifolds admit no nontrivial continuous circle action.
They fail the abelian fundamental-group hypothesis: the mapping-torus
group contains the nontrivial centerless group \(\pi_1(V)\), hence a
nonabelian subgroup \cite[Section~2, Step~4, manuscript p.~8]{CWY}.
Taking the product in their construction retains that subgroup.
They are not excluded merely by an effective-versus-free distinction.

The Wu manifold, its characteristic number and the circle action used here
are classical. This note assembles them to document the answer to the
printed 2010 question. An extensive but bounded priority audit on
16 October 2026 verified no complete predecessor meeting all its
hypotheses in the inspected statements. Some older books, theses,
journal versions and citation trails remain unread.
Historical priority remains unresolved; no firstness, new-ingredient,
certified-novelty or current worldwide-open claim is made.
Exact versions, read scopes and gaps are recorded in the companion archive.

The archive preserves the original candidate, independently positioned
mathematical and priority reviews, source identities, exact finite
weight/characteristic-algebra checks and reproduction instructions.
Those computations support the written arguments and do not formally
verify topology or establish an infinite quantifier by extrapolation.
Full third-party source bodies are not redistributed.

AI tools were used extensively in developing and checking the argument,
literature investigation, drafting, computation, reproduction and
adversarial review. This is an unrefereed preprint without external human
peer review or proof-assistant certification. The note and authored
companion material are licensed under CC BY 4.0.

\begin{thebibliography}{9}
\raggedright
\bibitem{H}
B.~Hanke, Homotopy Euler characteristic and the stable free rank of symmetry,
in \emph{Cohomology of Finite Groups: Interactions and Applications},
Oberwolfach Reports \textbf{7} (2010), report~32, pp.~1912--1915.
\href{https://doi.org/10.4171/OWR/2010/32}{doi:10.4171/OWR/2010/32}.
The report's publisher publication date is 2 March 2011.

\bibitem{DY}
A.~Debray and M.~Yu,
What bordism-theoretic anomaly cancellation can do for U,
\emph{Commun. Math. Phys.} \textbf{405} (2024), article~154.
\href{https://doi.org/10.1007/s00220-024-04937-4}{doi:10.1007/s00220-024-04937-4}.
Specific locators refer to the author-hosted arXiv v2 (header 22 July 2025;
internal Date line 23 July 2025), pp.~20--22: \url{https://adebray.github.io/papers/U_duality_Bordism.pdf}.

\bibitem{Y}
D.~Yeroshkin, Orbifold biquotients of \(SU(3)\),
\emph{Differential Geom. Appl.} \textbf{42} (2015), 54--76.
\href{https://doi.org/10.1016/j.difgeo.2015.07.003}{doi:10.1016/j.difgeo.2015.07.003}.
Proposition~4.2 is on pp.~6--7 of
\href{https://arxiv.org/abs/1401.7565v2}{arXiv:1401.7565v2} (9 April 2014).

\bibitem{KL}
N.~J.~Kuhn and C.~J.~R.~Lloyd,
Chromatic fixed point theory and the Balmer spectrum for extraspecial \(2\)-groups,
\href{https://arxiv.org/abs/2008.00330v1}{arXiv:2008.00330v1}
(1 August 2020), Theorem~9.1, p.~27;
\href{https://arxiv.org/abs/2008.00330v2}{v2}
(27 September 2022), Example~2.25, p.~12.

\bibitem{M}
I.~Mundet i Riera, Discrete degree of symmetry of manifolds,
\emph{Transformation Groups} \textbf{31} (2026), 801--838;
published online 19 April 2024.
\href{https://doi.org/10.1007/s00031-024-09858-z}{doi:10.1007/s00031-024-09858-z}.
The first preprint is
\href{https://arxiv.org/abs/2112.05599v1}{arXiv:2112.05599v1}
(10 December 2021).

\bibitem{CWY}
S.~Cappell, S.~Weinberger and M.~Yan,
Closed aspherical manifolds with center,
\emph{J. Topol.} \textbf{6} (2013), 1009--1018.
\href{https://doi.org/10.1112/jtopol/jtt023}{doi:10.1112/jtopol/jtt023}.
The manuscript locator refers to
\href{https://arxiv.org/abs/1108.2321v3}{arXiv:1108.2321v3}
(14 April 2013).
\end{thebibliography}
\end{document}
